% ==========================================
% PLANTILLA LATEX — DOCUMENTOS ACADÉMICOS
% ==========================================
%
% CÓMO USAR ESTA PLANTILLA:
%   1. Elige la fuente en  config/proyecto.tex
%   2. Edita  capitulos/portada.tex   con los datos de tu proyecto
%   3. Crea tus capítulos en  capitulos/
%   4. Descomenta los \input que necesites abajo
%   5. Compila:  latexmk -pdf -outdir=build main.tex
%
% ESTRUCTURA DE CARPETAS:
%   config/       → Configuración técnica (fuente, tamaños)
%   common/       → Archivos estáticos de la plantilla
%   capitulos/    → Portada y tus capítulos (editar aquí)
%   imagenes/     → Imágenes (logos en imagenes/logos/)
%   build/        → Salida de compilación (se genera sola, no se sube)
%
% TIP: Si algo no compila, limpia build/ y reintenta:
%      latexmk -C -outdir=build ; latexmk -pdf -outdir=build main.tex
%
% ==========================================

\documentclass[11pt, letterpaper]{report}

% --- Configuración técnica (fuente, tamaños) ---
% ==========================================
% CONFIGURACIÓN TÉCNICA
% ==========================================
%
% Los datos del proyecto (título, autores, etc.)
% se editan en: capitulos/portada.tex
%
% ==========================================

% --- Fuente principal ---
% "Alegreya" → fuente serif con carácter (XeLaTeX)
% "TimesNewRoman" → fuente clásica Times New Roman (XeLaTeX)
\newcommand{\fuentePrincipal}{Alegreya}

% --- Tamaños de fuente ---
\newcommand{\tamanoTextoBase}{11pt}
\newcommand{\tamanoTituloUno}{16pt}
\newcommand{\tamanoTituloDos}{14pt}
\newcommand{\tamanoTituloTres}{12pt}


% --- Paquetes y estilos (no tocar) ---
% ==========================================
% CONFIGURACIÓN DE LA PLANTILLA
% ==========================================
%
% ADVERTENCIA: No edites este archivo a menos que
% sepas lo que haces. La mayoría de las
% personalizaciones se hacen desde:
%   - capitulos/portada.tex  (datos del proyecto)
%   - config/proyecto.tex    (fuente, tamaños)
%
% ==========================================

% ==========================================
% PAQUETES ESENCIALES
% ==========================================

% --- Tipografía (XeLaTeX) ---
\usepackage{fontspec}

% Seleccionar fuente según \fuentePrincipal
\makeatletter
\newcommand{\@checkfont}{}
\let\@checkfont\fuentePrincipal
\newcommand{\@Alegreya}{Alegreya}
\newif\ifalegreya
\ifx\@checkfont\@Alegreya
  \alegreyatrue
\else
  \alegreyafalse
\fi
\makeatother

\ifalegreya
  \setmainfont{Alegreya}[
    SmallCapsFont={Alegreya SC},
    Numbers=OldStyle
  ]
  \newfontfamily{\AlegreyaSC}{Alegreya SC}
\else
  \setmainfont{Times New Roman}
  \newcommand{\AlegreyaSC}{\relax}
\fi

% Pasar opciones a xcolor
\PassOptionsToPackage{table,xcdraw}{xcolor}

% --- Idioma ---
\usepackage{csquotes}
\usepackage[spanish,es-tabla]{babel}

% --- Márgenes ---
\usepackage[top=2.5cm, bottom=2.5cm, left=3cm, right=2.5cm]{geometry}

% --- Gráficos e imágenes ---
\usepackage{graphicx}
\usepackage{adjustbox}
\usepackage{float}
\usepackage{caption}
\usepackage{subcaption}

\graphicspath{{imagenes/}}

% --- Colores ---
\usepackage{xcolor}

% --- Tablas ---
\usepackage{booktabs}
\usepackage{longtable}
\usepackage{multirow}

% --- Enlaces y referencias ---
\usepackage[hidelinks]{hyperref}
\usepackage{url}
\usepackage{seqsplit}

% --- Código fuente ---
\usepackage{listings}
\renewcommand{\lstlistingname}{Listado}
\lstset{
  inputencoding=utf8,
  extendedchars=true,
  literate={á}{{\'a}}1 {é}{{\'e}}1 {í}{{\'i}}1 {ó}{{\'o}}1 {ú}{{\'u}}1
  {Á}{{\'A}}1 {É}{{\'E}}1 {Í}{{\'I}}1 {Ó}{{\'O}}1 {Ú}{{\'U}}1
  {ñ}{{\~n}}1 {Ñ}{{\~N}}1 {¿}{{?`}}1 {¡}{{!`}}1,
  breaklines=true,
  breakatwhitespace=true,
  columns=fullflexible,
  keepspaces=true
}

% --- Encabezados y pies ---
\usepackage{fancyhdr}

% --- Espacio entre párrafos ---
\usepackage{parskip}

% --- Formato de títulos ---
\usepackage{titlesec}

\titleformat{\chapter}[display]
  {\normalfont\bfseries}{\chaptertitlename\ \thechapter}{20pt}
  {\fontsize{\tamanoTituloUno}{\dimexpr\tamanoTituloUno*12/10\relax}\selectfont}

\ifalegreya
  \titleformat{\section}
    {\AlegreyaSC\normalfont\bfseries}{\thesection}{1em}
    {\fontsize{\tamanoTituloDos}{\dimexpr\tamanoTituloDos*12/10\relax}\selectfont}
  \titleformat{\subsection}
    {\AlegreyaSC\normalfont\bfseries}{\thesubsection}{1em}
    {\fontsize{\tamanoTituloTres}{\dimexpr\tamanoTituloTres*12/10\relax}\selectfont}
\else
  \titleformat{\section}
    {\normalfont\bfseries}{\thesection}{1em}
    {\fontsize{\tamanoTituloDos}{\dimexpr\tamanoTituloDos*12/10\relax}\selectfont}
  \titleformat{\subsection}
    {\normalfont\bfseries}{\thesubsection}{1em}
    {\fontsize{\tamanoTituloTres}{\dimexpr\tamanoTituloTres*12/10\relax}\selectfont}
\fi

% --- Listas ---
\usepackage{enumitem}

% --- Matemáticas ---
\usepackage{amsmath}
\usepackage{amsfonts}
\usepackage{amssymb}

% --- Apéndices ---
\usepackage[titletoc]{appendix}

% --- Bibliografía ---
\usepackage[backend=biber, style=ieee, sorting=none]{biblatex}
\addbibresource{referencias.bib}

% ==========================================
% CONFIGURACIÓN DE ESTILOS
% ==========================================

% --- Colores para código ---
\definecolor{codegreen}{rgb}{0,0.6,0}
\definecolor{codegray}{rgb}{0.5,0.5,0.5}
\definecolor{codepurple}{rgb}{0.58,0,0.82}
\definecolor{backcolour}{rgb}{0.95,0.95,0.92}

% --- Estilo de listings ---
\lstdefinestyle{mystyle}{
    backgroundcolor=\color{backcolour},
    commentstyle=\color{codegreen},
    keywordstyle=\color{magenta},
    numberstyle=\tiny\color{codegray},
    stringstyle=\color{codepurple},
    basicstyle=\ttfamily\footnotesize,
    breakatwhitespace=false,
    breaklines=true,
    captionpos=b,
    keepspaces=true,
    numbers=left,
    numbersep=5pt,
    showspaces=false,
    showstringspaces=false,
    showtabs=false,
    tabsize=2
}
\lstset{style=mystyle}

% --- Encabezados ---
\setlength{\headheight}{14.2pt}
\pagestyle{fancy}
\fancyhf{}
\fancyhead[L]{\leftmark}
\fancyhead[R]{\thepage}
\renewcommand{\headrulewidth}{0.5pt}

% ==========================================
% COMANDOS ÚTILES
% ==========================================

\newcommand{\comillas}[1]{``#1''}
\newcommand{\codigo}[1]{\colorbox{backcolour}{\texttt{#1}}}
\newcommand{\nota}[1]{\marginpar{\footnotesize\textit{#1}}}
\newcommand{\propia}{Fuente: Elaboración propia.}


% ==========================================
\begin{document}
% ==========================================

% --- Portada ---
\pagenumbering{gobble}
% ==========================================
% PORTADA — EDITAR PARA CADA PROYECTO
% ==========================================
%
% Edita los datos marcados con (*).
% La fuente se elige en: config/proyecto.tex
% ==========================================

% ==========================================
% DATOS DEL PROYECTO (*)
% ==========================================

\newcommand{\titulo}{Título del proyecto}
\newcommand{\subtitulo}{Subtítulo o descripción corta}
\newcommand{\materia}{Nombre de la materia}
\newcommand{\profesor}{Nombre del profesor}
\newcommand{\fechaEntrega}{\today}

\newcommand{\institucion}{Instituto Politécnico Nacional}
\newcommand{\escuela}{Escuela Superior de Cómputo}
\newcommand{\siglasEscuela}{ESCOM}
\newcommand{\carrera}{Ingeniería en Inteligencia Artificial}
\newcommand{\grupo}{XBV1}
\newcommand{\lugar}{Ciudad de México}

\newcommand{\listaAutores}{
    Apellido1 Nombre1 \\
    Apellido2 Nombre2
}

\newcommand{\logoIzquierdo}{imagenes/logos/ipn_logo.png}
\newcommand{\logoDerecho}{imagenes/logos/escom_logo.png}

% ==========================================
% DISEÑO DE LA PORTADA
% ==========================================

\begin{titlepage}
    \newgeometry{left=2cm, right=2cm, top=2.5cm, bottom=2.5cm}
    \centering

    \begin{minipage}{0.18\textwidth}
        \centering
        \ifx\logoIzquierdo\empty\else
            \includegraphics[width=0.7\textwidth]{\logoIzquierdo}
        \fi
    \end{minipage}%
    \begin{minipage}{0.64\textwidth}
        \centering
        {\AlegreyaSC\fontsize{21pt}{28pt}\selectfont\bfseries \institucion}\\[0.4cm]
        {\AlegreyaSC\fontsize{19pt}{26pt}\selectfont\bfseries \escuela}\\[0.35cm]
        {\AlegreyaSC\fontsize{17pt}{19pt}\selectfont\bfseries \siglasEscuela}
    \end{minipage}%
    \begin{minipage}{0.18\textwidth}
        \centering
        \ifx\logoDerecho\empty\else
            \includegraphics[width=\textwidth]{\logoDerecho}
        \fi
    \end{minipage}

    \vspace{1.3cm}

    {\AlegreyaSC\LARGE \carrera}\\[1cm]
    {\AlegreyaSC\Large \materia}\\[1.5cm]

    \rule{3cm}{2pt}\\[1cm]

    {\fontsize{28pt}{34pt}\selectfont\itshape \titulo}\\[0.5cm]

    \ifx\subtitulo\empty\else
        {\Large\itshape \subtitulo}\\[0.5cm]
    \fi

    \vspace{1.5cm}

    {\Large\textbf{INTEGRANTES}}\\[0.8cm]
    \begin{center}
        \large
        \begin{tabular}{c}
            \listaAutores
        \end{tabular}
    \end{center}

    \ifx\profesor\empty\else
        \vspace{1cm}
        {\Large\textbf{Profesor:} \profesor}\\[0.5cm]
    \fi

    \vspace{\fill}

    {\large\textbf{Grupo:} \grupo}\\[0.5cm]
    {\large \lugar\ a \fechaEntrega}

    \restoregeometry
\end{titlepage}


% --- Resumen ---
% \input{capitulos/00_resumen}

% --- Índices ---
% ==========================================
% ÍNDICES — ARCHIVO ESTÁTICO
% ==========================================
%
% Se generan automáticamente a partir del contenido.
% ==========================================

\pagenumbering{roman}
\setcounter{page}{1}

\tableofcontents
\newpage

\listoffigures
\newpage

\listoftables
\newpage


% --- Contenido principal ---
\pagenumbering{arabic}

% TIP: Descomenta y agrega los capítulos que necesites.
% \input{capitulos/introduccion}
% \input{capitulos/marco_teorico}
% \input{capitulos/desarrollo}
% \input{capitulos/conclusiones}

% Capítulo de ejemplo (descomentar para ver la referencia de LaTeX)
% % ==========================================
% CAPÍTULO DE EJEMPLO
% ==========================================
% Este capítulo muestra los patrones más comunes de LaTeX.
% Úsalo como referencia y elimínalo cuando ya no lo necesites.
%
% TIP: Cada capítulo es un archivo independiente en capitulos/
%      y se incluye desde main.tex con \input{capitulos/nombre}
%
% TIP: Nunca pongas \begin{document} ni \end{document} aquí.
%      Eso sólo va en main.tex.
% ==========================================

\chapter{Capítulo de ejemplo}
\label{cap:ejemplo}

% TIP: \label{} te permite referenciar este capítulo desde cualquier parte:
%      "Como se menciona en el Capítulo \ref{cap:ejemplo}..."

\section{Texto básico}

Este es un párrafo normal. En \LaTeX{}, los párrafos se separan con una línea en blanco.
No uses \texttt{\textbackslash\textbackslash} para saltos de línea en texto normal,
eso es sólo para tablas y ecuaciones.

% TIP: Para texto en negritas usa \textbf{}, para itálicas \textit{},
%      para código en línea \texttt{} o el comando \codigo{} definido en la plantilla.

Este texto tiene \textbf{negritas}, \textit{itálicas} y \codigo{código en línea}.
También puedes usar las \comillas{comillas tipográficas} definidas en la plantilla.


\section{Figuras}
\label{sec:figuras}

% TIP: Coloca tus imágenes en la carpeta imagenes/
%      LaTeX las encuentra automáticamente gracias a \graphicspath en proyecto.tex
%
% TIP: Usa [H] (requiere paquete float) para forzar la posición exacta.
%      Usa [htbp] para dejar que LaTeX elija la mejor posición.
%      Preferir [htbp] produce documentos mejor maquetados.
%
% TIP: SIEMPRE pon \label DESPUÉS de \caption, nunca antes.
%      De lo contrario, las referencias cruzadas apuntarán al lugar incorrecto.
%
% TIP: No dejes líneas en blanco entre \caption y \label.

% --- Ejemplo de figura simple ---
% \begin{figure}[htbp]
%     \centering
%     \includegraphics[width=0.7\textwidth]{nombre_imagen.png}
%     \caption{Descripción clara de la imagen}
%     \label{fig:mi_figura}
% \end{figure}

% Para referenciar: "Como se muestra en la Figura \ref{fig:mi_figura}..."


\section{Tablas}
\label{sec:tablas}

% TIP: Usa booktabs (\toprule, \midrule, \bottomrule) en lugar de \hline.
%      Se ve mucho más profesional.
%
% TIP: No uses líneas verticales (|) en tablas académicas.
%
% TIP: Para tablas largas que cruzan páginas, usa longtable.
%      Para tablas cortas, usa table + tabular.

% --- Ejemplo de tabla corta ---
\begin{table}[htbp]
\centering
\begin{tabular}{@{}llr@{}}
\toprule
\textbf{Elemento} & \textbf{Categoría} & \textbf{Valor} \\ \midrule
Primer elemento    & Tipo A             & 100             \\
Segundo elemento   & Tipo B             & 200             \\
Tercer elemento    & Tipo A             & 150             \\ \midrule
\textbf{Total}     &                    & \textbf{450}    \\ \bottomrule
\end{tabular}
\caption{Ejemplo de tabla con booktabs}
\label{tab:ejemplo}
\end{table}

% --- Ejemplo de tabla larga (multi-página) ---
% \begin{longtable}{@{}p{2cm}p{8cm}p{3cm}@{}}
% \toprule
% \textbf{ID} & \textbf{Descripción} & \textbf{Estado} \\ \midrule
% \endhead  % <-- Esto repite el encabezado en cada página
% Fila 1 & Descripción de la fila 1 & Activo \\
% Fila 2 & Descripción de la fila 2 & Inactivo \\ \bottomrule
% \caption{Ejemplo de tabla larga}
% \label{tab:ejemplo_largo}
% \end{longtable}


\section{Ecuaciones}
\label{sec:ecuaciones}

% TIP: Ecuación en línea: $E = mc^2$
% TIP: Ecuación centrada sin número: \[ E = mc^2 \]
% TIP: Ecuación centrada con número (referenciable): usar equation

La famosa ecuación de Einstein es $E = mc^2$.

% --- Ejemplo de ecuación numerada ---
\begin{equation}
    \sum_{i=1}^{n} i = \frac{n(n+1)}{2}
    \label{eq:suma_gauss}
\end{equation}

La Ecuación \ref{eq:suma_gauss} muestra la suma de Gauss.


\section{Listas}
\label{sec:listas}

% TIP: itemize = viñetas, enumerate = numerada, description = con título

% --- Lista con viñetas ---
\begin{itemize}
    \item Primer punto
    \item Segundo punto
    \item Tercer punto con sublista:
    \begin{itemize}
        \item Subpunto A
        \item Subpunto B
    \end{itemize}
\end{itemize}

% --- Lista numerada ---
\begin{enumerate}
    \item Primer paso
    \item Segundo paso
    \item Tercer paso
\end{enumerate}


\section{Código fuente}
\label{sec:codigo}

% TIP: Usa lstlisting para bloques de código.
%      El estilo 'mystyle' ya está configurado en configuracion.tex.
%
% TIP: Para lenguajes soportados, usa language=Python, language=Java, etc.
%      Esto activa el resaltado de sintaxis automático.
%
% TIP: Si necesitas resaltado más avanzado, activa minted en configuracion.tex
%      (requiere Python + Pygments instalado).

\begin{lstlisting}[language=Python, caption={Ejemplo de código Python}, label={lst:ejemplo_python}]
def fibonacci(n):
    """Calcula el n-esimo numero de Fibonacci."""
    if n <= 1:
        return n
    return fibonacci(n - 1) + fibonacci(n - 2)

# Ejemplo de uso
for i in range(10):
    print(f"F({i}) = {fibonacci(i)}")
\end{lstlisting}

El Listado \ref{lst:ejemplo_python} muestra un ejemplo de código con resaltado de sintaxis. 


\section{Citas bibliográficas}
\label{sec:citas}

% TIP: Las referencias se definen en referencias.bib
%      y se citan con \cite{clave} o \textcite{clave}.
%
% TIP: Ejemplo de entrada en referencias.bib:
%
% @book{knuth1997,
%     author    = {Donald E. Knuth},
%     title     = {The Art of Computer Programming},
%     publisher = {Addison-Wesley},
%     year      = {1997},
%     volume    = {1}
% }
%
% Luego en el texto: "Según Knuth \cite{knuth1997}..."
%
% TIP: Para que las citas aparezcan, descomenta % ==========================================
% BIBLIOGRAFÍA — ARCHIVO ESTÁTICO
% ==========================================
%
% No edites este archivo directamente.
% Agrega tus referencias en: referencias.bib
%
% TIP: Compilación con bibliografía:
%      xelatex → biber → xelatex → xelatex
%      O simplemente: latexmk -pdf -outdir=build main.tex
% ==========================================

\begingroup
\sloppy
\printbibliography[heading=bibintoc, title={Referencias}]
\endgroup

%      en main.tex y compila con: pdflatex → biber → pdflatex → pdflatex


\section{Referencias cruzadas}
\label{sec:referencias}

% TIP: Puedes referenciar cualquier elemento con \label y \ref:
%      Capítulos: Capítulo \ref{cap:ejemplo}
%      Secciones: Sección \ref{sec:figuras}
%      Figuras:   Figura \ref{fig:mi_figura}
%      Tablas:    Tabla \ref{tab:ejemplo}
%      Ecuaciones: Ecuación \ref{eq:suma_gauss}
%      Código:    Listado \ref{lst:ejemplo_python}
%
% TIP: Las referencias se actualizan al compilar DOS veces.
%      Si ves "??" en lugar de números, compila de nuevo.

Este capítulo (Capítulo \ref{cap:ejemplo}) contiene ejemplos de:
tablas (Tabla~\ref{tab:ejemplo}),
ecuaciones (Ecuación~\ref{eq:suma_gauss})
y código (Listado~\ref{lst:ejemplo_python}).

% TIP: Usa ~ (tilde) entre "Figura" y \ref para evitar que
%      se separen en líneas diferentes. Ej: Figura~\ref{fig:x}


% --- Bibliografía ---
% TIP: Si referencias.bib está vacío, comenta también \nocite{*} o biber fallará.
% \nocite{*}
% % ==========================================
% BIBLIOGRAFÍA — ARCHIVO ESTÁTICO
% ==========================================
%
% No edites este archivo directamente.
% Agrega tus referencias en: referencias.bib
%
% TIP: Compilación con bibliografía:
%      xelatex → biber → xelatex → xelatex
%      O simplemente: latexmk -pdf -outdir=build main.tex
% ==========================================

\begingroup
\sloppy
\printbibliography[heading=bibintoc, title={Referencias}]
\endgroup


% --- Anexos ---
% % ==========================================
% ANEXOS — ARCHIVO AGREGADOR
% ==========================================
%
% Este archivo es el punto de entrada de todos los anexos.
% No pongas \chapter{} aquí directamente (a menos que sea
% un anexo general). En su lugar, crea un anexo por capítulo
% y agrégalo con \input{}.
%
% TIP: Cada capítulo puede tener su propio anexo en:
%      capitulos/capitulo_X/anexos.tex
%
% TIP: Los anexos se numeran con letras (A, B, C...)
%      automáticamente según el orden de los \input{}.
% ==========================================

\appendix
\renewcommand{\appendixname}{Anexo}
\renewcommand{\appendixpagename}{Anexos}
\renewcommand{\appendixtocname}{Anexos}
\addappheadtotoc
\appendixpage

% --- Agrega aquí los anexos de cada capítulo ---
% Ejemplo:
% % ==========================================
% ANEXOS — ARCHIVO AGREGADOR
% ==========================================
%
% Este archivo es el punto de entrada de todos los anexos.
% No pongas \chapter{} aquí directamente (a menos que sea
% un anexo general). En su lugar, crea un anexo por capítulo
% y agrégalo con \input{}.
%
% TIP: Cada capítulo puede tener su propio anexo en:
%      capitulos/capitulo_X/anexos.tex
%
% TIP: Los anexos se numeran con letras (A, B, C...)
%      automáticamente según el orden de los \input{}.
% ==========================================

\appendix
\renewcommand{\appendixname}{Anexo}
\renewcommand{\appendixpagename}{Anexos}
\renewcommand{\appendixtocname}{Anexos}
\addappheadtotoc
\appendixpage

% --- Agrega aquí los anexos de cada capítulo ---
% Ejemplo:
% % ==========================================
% ANEXOS — ARCHIVO AGREGADOR
% ==========================================
%
% Este archivo es el punto de entrada de todos los anexos.
% No pongas \chapter{} aquí directamente (a menos que sea
% un anexo general). En su lugar, crea un anexo por capítulo
% y agrégalo con \input{}.
%
% TIP: Cada capítulo puede tener su propio anexo en:
%      capitulos/capitulo_X/anexos.tex
%
% TIP: Los anexos se numeran con letras (A, B, C...)
%      automáticamente según el orden de los \input{}.
% ==========================================

\appendix
\renewcommand{\appendixname}{Anexo}
\renewcommand{\appendixpagename}{Anexos}
\renewcommand{\appendixtocname}{Anexos}
\addappheadtotoc
\appendixpage

% --- Agrega aquí los anexos de cada capítulo ---
% Ejemplo:
% \input{capitulos/capitulo_1/anexos}
% \input{capitulos/capitulo_2/anexos}

% ==========================================
% REFERENCIA: Cómo incluir código fuente
% ==========================================
%
% Para proyectos pequeños, puedes poner los anexos
% directamente aquí. Para proyectos grandes, crea un
% archivo por capítulo (ver sección "Documentos grandes"
% en el README.md).
%
% A continuación se muestran las formas más comunes
% de incluir código con listings.
%
% ==========================================
%
% --- FORMA 1: Desde archivo externo (recomendada) ---
% Crea una carpeta codigos/ o programa/ en la raíz.
%
% \chapter{Código fuente}
% \label{anexo:codigo}
%
% \lstinputlisting[language=C, style=mystyle,
%   caption={Descripción del código. \propia{}},
%   label={lst:etiqueta}]{codigos/mi_archivo.c}
%
%
% --- FORMA 2: Con verificación de existencia ---
% Si el archivo puede no estar presente, usa \IfFileExists
% para evitar errores de compilación.
%
% \IfFileExists{programa/mi_programa.c}{
%   \lstinputlisting[language=C, style=mystyle,
%     caption={Código completo. \propia{}},
%     label={lst:codigo}]{programa/mi_programa.c}
% }{
%   \noindent\textbf{Nota:} El archivo aún no está disponible.
% }
%
%
% --- FORMA 3: Código inline (fragmentos pequeños) ---
% \begin{lstlisting}[language=Python, style=mystyle,
%   caption={Fragmento de ejemplo},
%   label={lst:ejemplo}]
% def fibonacci(n):
%     if n <= 1:
%         return n
%     return fibonacci(n - 1) + fibonacci(n - 2)
% \end{lstlisting}
%
%
% --- FORMA 4: Makefile ---
% \lstinputlisting[language=make, style=mystyle,
%   caption={Makefile del proyecto. \propia{}},
%   label={lst:makefile}]{codigos/Makefile}
%
%
% --- FORMA 5: Archivos de cabecera (.h) ---
% \lstinputlisting[language=C, style=mystyle,
%   caption={Archivo de cabecera. \propia{}},
%   label={lst:header}]{programa/common.h}
%
%
% ==========================================
% BUENAS PRÁCTICAS
% ==========================================
%
% 1. Usa \lstinputlisting para código que compila.
%    Así el PDF siempre refleja el código real.
%
% 2. Usa \IfFileExists para que la compilación no
%    falle si falta un archivo.
%
% 3. Pon siempre caption y label para referenciar
%    los listados desde el texto:
%    "En el Listado~\ref{lst:etiqueta} se muestra..."
%
% 4. Usa \propia{} al final del caption para agregar
%    "Fuente: Elaboración propia."
%
% 5. Para proyectos con varios capítulos, crea un
%    anexos.tex por capítulo y agrégalo arriba con
%    \input{}.

% % ==========================================
% ANEXOS — ARCHIVO AGREGADOR
% ==========================================
%
% Este archivo es el punto de entrada de todos los anexos.
% No pongas \chapter{} aquí directamente (a menos que sea
% un anexo general). En su lugar, crea un anexo por capítulo
% y agrégalo con \input{}.
%
% TIP: Cada capítulo puede tener su propio anexo en:
%      capitulos/capitulo_X/anexos.tex
%
% TIP: Los anexos se numeran con letras (A, B, C...)
%      automáticamente según el orden de los \input{}.
% ==========================================

\appendix
\renewcommand{\appendixname}{Anexo}
\renewcommand{\appendixpagename}{Anexos}
\renewcommand{\appendixtocname}{Anexos}
\addappheadtotoc
\appendixpage

% --- Agrega aquí los anexos de cada capítulo ---
% Ejemplo:
% \input{capitulos/capitulo_1/anexos}
% \input{capitulos/capitulo_2/anexos}

% ==========================================
% REFERENCIA: Cómo incluir código fuente
% ==========================================
%
% Para proyectos pequeños, puedes poner los anexos
% directamente aquí. Para proyectos grandes, crea un
% archivo por capítulo (ver sección "Documentos grandes"
% en el README.md).
%
% A continuación se muestran las formas más comunes
% de incluir código con listings.
%
% ==========================================
%
% --- FORMA 1: Desde archivo externo (recomendada) ---
% Crea una carpeta codigos/ o programa/ en la raíz.
%
% \chapter{Código fuente}
% \label{anexo:codigo}
%
% \lstinputlisting[language=C, style=mystyle,
%   caption={Descripción del código. \propia{}},
%   label={lst:etiqueta}]{codigos/mi_archivo.c}
%
%
% --- FORMA 2: Con verificación de existencia ---
% Si el archivo puede no estar presente, usa \IfFileExists
% para evitar errores de compilación.
%
% \IfFileExists{programa/mi_programa.c}{
%   \lstinputlisting[language=C, style=mystyle,
%     caption={Código completo. \propia{}},
%     label={lst:codigo}]{programa/mi_programa.c}
% }{
%   \noindent\textbf{Nota:} El archivo aún no está disponible.
% }
%
%
% --- FORMA 3: Código inline (fragmentos pequeños) ---
% \begin{lstlisting}[language=Python, style=mystyle,
%   caption={Fragmento de ejemplo},
%   label={lst:ejemplo}]
% def fibonacci(n):
%     if n <= 1:
%         return n
%     return fibonacci(n - 1) + fibonacci(n - 2)
% \end{lstlisting}
%
%
% --- FORMA 4: Makefile ---
% \lstinputlisting[language=make, style=mystyle,
%   caption={Makefile del proyecto. \propia{}},
%   label={lst:makefile}]{codigos/Makefile}
%
%
% --- FORMA 5: Archivos de cabecera (.h) ---
% \lstinputlisting[language=C, style=mystyle,
%   caption={Archivo de cabecera. \propia{}},
%   label={lst:header}]{programa/common.h}
%
%
% ==========================================
% BUENAS PRÁCTICAS
% ==========================================
%
% 1. Usa \lstinputlisting para código que compila.
%    Así el PDF siempre refleja el código real.
%
% 2. Usa \IfFileExists para que la compilación no
%    falle si falta un archivo.
%
% 3. Pon siempre caption y label para referenciar
%    los listados desde el texto:
%    "En el Listado~\ref{lst:etiqueta} se muestra..."
%
% 4. Usa \propia{} al final del caption para agregar
%    "Fuente: Elaboración propia."
%
% 5. Para proyectos con varios capítulos, crea un
%    anexos.tex por capítulo y agrégalo arriba con
%    \input{}.


% ==========================================
% REFERENCIA: Cómo incluir código fuente
% ==========================================
%
% Para proyectos pequeños, puedes poner los anexos
% directamente aquí. Para proyectos grandes, crea un
% archivo por capítulo (ver sección "Documentos grandes"
% en el README.md).
%
% A continuación se muestran las formas más comunes
% de incluir código con listings.
%
% ==========================================
%
% --- FORMA 1: Desde archivo externo (recomendada) ---
% Crea una carpeta codigos/ o programa/ en la raíz.
%
% \chapter{Código fuente}
% \label{anexo:codigo}
%
% \lstinputlisting[language=C, style=mystyle,
%   caption={Descripción del código. \propia{}},
%   label={lst:etiqueta}]{codigos/mi_archivo.c}
%
%
% --- FORMA 2: Con verificación de existencia ---
% Si el archivo puede no estar presente, usa \IfFileExists
% para evitar errores de compilación.
%
% \IfFileExists{programa/mi_programa.c}{
%   \lstinputlisting[language=C, style=mystyle,
%     caption={Código completo. \propia{}},
%     label={lst:codigo}]{programa/mi_programa.c}
% }{
%   \noindent\textbf{Nota:} El archivo aún no está disponible.
% }
%
%
% --- FORMA 3: Código inline (fragmentos pequeños) ---
% \begin{lstlisting}[language=Python, style=mystyle,
%   caption={Fragmento de ejemplo},
%   label={lst:ejemplo}]
% def fibonacci(n):
%     if n <= 1:
%         return n
%     return fibonacci(n - 1) + fibonacci(n - 2)
% \end{lstlisting}
%
%
% --- FORMA 4: Makefile ---
% \lstinputlisting[language=make, style=mystyle,
%   caption={Makefile del proyecto. \propia{}},
%   label={lst:makefile}]{codigos/Makefile}
%
%
% --- FORMA 5: Archivos de cabecera (.h) ---
% \lstinputlisting[language=C, style=mystyle,
%   caption={Archivo de cabecera. \propia{}},
%   label={lst:header}]{programa/common.h}
%
%
% ==========================================
% BUENAS PRÁCTICAS
% ==========================================
%
% 1. Usa \lstinputlisting para código que compila.
%    Así el PDF siempre refleja el código real.
%
% 2. Usa \IfFileExists para que la compilación no
%    falle si falta un archivo.
%
% 3. Pon siempre caption y label para referenciar
%    los listados desde el texto:
%    "En el Listado~\ref{lst:etiqueta} se muestra..."
%
% 4. Usa \propia{} al final del caption para agregar
%    "Fuente: Elaboración propia."
%
% 5. Para proyectos con varios capítulos, crea un
%    anexos.tex por capítulo y agrégalo arriba con
%    \input{}.

% % ==========================================
% ANEXOS — ARCHIVO AGREGADOR
% ==========================================
%
% Este archivo es el punto de entrada de todos los anexos.
% No pongas \chapter{} aquí directamente (a menos que sea
% un anexo general). En su lugar, crea un anexo por capítulo
% y agrégalo con \input{}.
%
% TIP: Cada capítulo puede tener su propio anexo en:
%      capitulos/capitulo_X/anexos.tex
%
% TIP: Los anexos se numeran con letras (A, B, C...)
%      automáticamente según el orden de los \input{}.
% ==========================================

\appendix
\renewcommand{\appendixname}{Anexo}
\renewcommand{\appendixpagename}{Anexos}
\renewcommand{\appendixtocname}{Anexos}
\addappheadtotoc
\appendixpage

% --- Agrega aquí los anexos de cada capítulo ---
% Ejemplo:
% % ==========================================
% ANEXOS — ARCHIVO AGREGADOR
% ==========================================
%
% Este archivo es el punto de entrada de todos los anexos.
% No pongas \chapter{} aquí directamente (a menos que sea
% un anexo general). En su lugar, crea un anexo por capítulo
% y agrégalo con \input{}.
%
% TIP: Cada capítulo puede tener su propio anexo en:
%      capitulos/capitulo_X/anexos.tex
%
% TIP: Los anexos se numeran con letras (A, B, C...)
%      automáticamente según el orden de los \input{}.
% ==========================================

\appendix
\renewcommand{\appendixname}{Anexo}
\renewcommand{\appendixpagename}{Anexos}
\renewcommand{\appendixtocname}{Anexos}
\addappheadtotoc
\appendixpage

% --- Agrega aquí los anexos de cada capítulo ---
% Ejemplo:
% \input{capitulos/capitulo_1/anexos}
% \input{capitulos/capitulo_2/anexos}

% ==========================================
% REFERENCIA: Cómo incluir código fuente
% ==========================================
%
% Para proyectos pequeños, puedes poner los anexos
% directamente aquí. Para proyectos grandes, crea un
% archivo por capítulo (ver sección "Documentos grandes"
% en el README.md).
%
% A continuación se muestran las formas más comunes
% de incluir código con listings.
%
% ==========================================
%
% --- FORMA 1: Desde archivo externo (recomendada) ---
% Crea una carpeta codigos/ o programa/ en la raíz.
%
% \chapter{Código fuente}
% \label{anexo:codigo}
%
% \lstinputlisting[language=C, style=mystyle,
%   caption={Descripción del código. \propia{}},
%   label={lst:etiqueta}]{codigos/mi_archivo.c}
%
%
% --- FORMA 2: Con verificación de existencia ---
% Si el archivo puede no estar presente, usa \IfFileExists
% para evitar errores de compilación.
%
% \IfFileExists{programa/mi_programa.c}{
%   \lstinputlisting[language=C, style=mystyle,
%     caption={Código completo. \propia{}},
%     label={lst:codigo}]{programa/mi_programa.c}
% }{
%   \noindent\textbf{Nota:} El archivo aún no está disponible.
% }
%
%
% --- FORMA 3: Código inline (fragmentos pequeños) ---
% \begin{lstlisting}[language=Python, style=mystyle,
%   caption={Fragmento de ejemplo},
%   label={lst:ejemplo}]
% def fibonacci(n):
%     if n <= 1:
%         return n
%     return fibonacci(n - 1) + fibonacci(n - 2)
% \end{lstlisting}
%
%
% --- FORMA 4: Makefile ---
% \lstinputlisting[language=make, style=mystyle,
%   caption={Makefile del proyecto. \propia{}},
%   label={lst:makefile}]{codigos/Makefile}
%
%
% --- FORMA 5: Archivos de cabecera (.h) ---
% \lstinputlisting[language=C, style=mystyle,
%   caption={Archivo de cabecera. \propia{}},
%   label={lst:header}]{programa/common.h}
%
%
% ==========================================
% BUENAS PRÁCTICAS
% ==========================================
%
% 1. Usa \lstinputlisting para código que compila.
%    Así el PDF siempre refleja el código real.
%
% 2. Usa \IfFileExists para que la compilación no
%    falle si falta un archivo.
%
% 3. Pon siempre caption y label para referenciar
%    los listados desde el texto:
%    "En el Listado~\ref{lst:etiqueta} se muestra..."
%
% 4. Usa \propia{} al final del caption para agregar
%    "Fuente: Elaboración propia."
%
% 5. Para proyectos con varios capítulos, crea un
%    anexos.tex por capítulo y agrégalo arriba con
%    \input{}.

% % ==========================================
% ANEXOS — ARCHIVO AGREGADOR
% ==========================================
%
% Este archivo es el punto de entrada de todos los anexos.
% No pongas \chapter{} aquí directamente (a menos que sea
% un anexo general). En su lugar, crea un anexo por capítulo
% y agrégalo con \input{}.
%
% TIP: Cada capítulo puede tener su propio anexo en:
%      capitulos/capitulo_X/anexos.tex
%
% TIP: Los anexos se numeran con letras (A, B, C...)
%      automáticamente según el orden de los \input{}.
% ==========================================

\appendix
\renewcommand{\appendixname}{Anexo}
\renewcommand{\appendixpagename}{Anexos}
\renewcommand{\appendixtocname}{Anexos}
\addappheadtotoc
\appendixpage

% --- Agrega aquí los anexos de cada capítulo ---
% Ejemplo:
% \input{capitulos/capitulo_1/anexos}
% \input{capitulos/capitulo_2/anexos}

% ==========================================
% REFERENCIA: Cómo incluir código fuente
% ==========================================
%
% Para proyectos pequeños, puedes poner los anexos
% directamente aquí. Para proyectos grandes, crea un
% archivo por capítulo (ver sección "Documentos grandes"
% en el README.md).
%
% A continuación se muestran las formas más comunes
% de incluir código con listings.
%
% ==========================================
%
% --- FORMA 1: Desde archivo externo (recomendada) ---
% Crea una carpeta codigos/ o programa/ en la raíz.
%
% \chapter{Código fuente}
% \label{anexo:codigo}
%
% \lstinputlisting[language=C, style=mystyle,
%   caption={Descripción del código. \propia{}},
%   label={lst:etiqueta}]{codigos/mi_archivo.c}
%
%
% --- FORMA 2: Con verificación de existencia ---
% Si el archivo puede no estar presente, usa \IfFileExists
% para evitar errores de compilación.
%
% \IfFileExists{programa/mi_programa.c}{
%   \lstinputlisting[language=C, style=mystyle,
%     caption={Código completo. \propia{}},
%     label={lst:codigo}]{programa/mi_programa.c}
% }{
%   \noindent\textbf{Nota:} El archivo aún no está disponible.
% }
%
%
% --- FORMA 3: Código inline (fragmentos pequeños) ---
% \begin{lstlisting}[language=Python, style=mystyle,
%   caption={Fragmento de ejemplo},
%   label={lst:ejemplo}]
% def fibonacci(n):
%     if n <= 1:
%         return n
%     return fibonacci(n - 1) + fibonacci(n - 2)
% \end{lstlisting}
%
%
% --- FORMA 4: Makefile ---
% \lstinputlisting[language=make, style=mystyle,
%   caption={Makefile del proyecto. \propia{}},
%   label={lst:makefile}]{codigos/Makefile}
%
%
% --- FORMA 5: Archivos de cabecera (.h) ---
% \lstinputlisting[language=C, style=mystyle,
%   caption={Archivo de cabecera. \propia{}},
%   label={lst:header}]{programa/common.h}
%
%
% ==========================================
% BUENAS PRÁCTICAS
% ==========================================
%
% 1. Usa \lstinputlisting para código que compila.
%    Así el PDF siempre refleja el código real.
%
% 2. Usa \IfFileExists para que la compilación no
%    falle si falta un archivo.
%
% 3. Pon siempre caption y label para referenciar
%    los listados desde el texto:
%    "En el Listado~\ref{lst:etiqueta} se muestra..."
%
% 4. Usa \propia{} al final del caption para agregar
%    "Fuente: Elaboración propia."
%
% 5. Para proyectos con varios capítulos, crea un
%    anexos.tex por capítulo y agrégalo arriba con
%    \input{}.


% ==========================================
% REFERENCIA: Cómo incluir código fuente
% ==========================================
%
% Para proyectos pequeños, puedes poner los anexos
% directamente aquí. Para proyectos grandes, crea un
% archivo por capítulo (ver sección "Documentos grandes"
% en el README.md).
%
% A continuación se muestran las formas más comunes
% de incluir código con listings.
%
% ==========================================
%
% --- FORMA 1: Desde archivo externo (recomendada) ---
% Crea una carpeta codigos/ o programa/ en la raíz.
%
% \chapter{Código fuente}
% \label{anexo:codigo}
%
% \lstinputlisting[language=C, style=mystyle,
%   caption={Descripción del código. \propia{}},
%   label={lst:etiqueta}]{codigos/mi_archivo.c}
%
%
% --- FORMA 2: Con verificación de existencia ---
% Si el archivo puede no estar presente, usa \IfFileExists
% para evitar errores de compilación.
%
% \IfFileExists{programa/mi_programa.c}{
%   \lstinputlisting[language=C, style=mystyle,
%     caption={Código completo. \propia{}},
%     label={lst:codigo}]{programa/mi_programa.c}
% }{
%   \noindent\textbf{Nota:} El archivo aún no está disponible.
% }
%
%
% --- FORMA 3: Código inline (fragmentos pequeños) ---
% \begin{lstlisting}[language=Python, style=mystyle,
%   caption={Fragmento de ejemplo},
%   label={lst:ejemplo}]
% def fibonacci(n):
%     if n <= 1:
%         return n
%     return fibonacci(n - 1) + fibonacci(n - 2)
% \end{lstlisting}
%
%
% --- FORMA 4: Makefile ---
% \lstinputlisting[language=make, style=mystyle,
%   caption={Makefile del proyecto. \propia{}},
%   label={lst:makefile}]{codigos/Makefile}
%
%
% --- FORMA 5: Archivos de cabecera (.h) ---
% \lstinputlisting[language=C, style=mystyle,
%   caption={Archivo de cabecera. \propia{}},
%   label={lst:header}]{programa/common.h}
%
%
% ==========================================
% BUENAS PRÁCTICAS
% ==========================================
%
% 1. Usa \lstinputlisting para código que compila.
%    Así el PDF siempre refleja el código real.
%
% 2. Usa \IfFileExists para que la compilación no
%    falle si falta un archivo.
%
% 3. Pon siempre caption y label para referenciar
%    los listados desde el texto:
%    "En el Listado~\ref{lst:etiqueta} se muestra..."
%
% 4. Usa \propia{} al final del caption para agregar
%    "Fuente: Elaboración propia."
%
% 5. Para proyectos con varios capítulos, crea un
%    anexos.tex por capítulo y agrégalo arriba con
%    \input{}.


% ==========================================
% REFERENCIA: Cómo incluir código fuente
% ==========================================
%
% Para proyectos pequeños, puedes poner los anexos
% directamente aquí. Para proyectos grandes, crea un
% archivo por capítulo (ver sección "Documentos grandes"
% en el README.md).
%
% A continuación se muestran las formas más comunes
% de incluir código con listings.
%
% ==========================================
%
% --- FORMA 1: Desde archivo externo (recomendada) ---
% Crea una carpeta codigos/ o programa/ en la raíz.
%
% \chapter{Código fuente}
% \label{anexo:codigo}
%
% \lstinputlisting[language=C, style=mystyle,
%   caption={Descripción del código. \propia{}},
%   label={lst:etiqueta}]{codigos/mi_archivo.c}
%
%
% --- FORMA 2: Con verificación de existencia ---
% Si el archivo puede no estar presente, usa \IfFileExists
% para evitar errores de compilación.
%
% \IfFileExists{programa/mi_programa.c}{
%   \lstinputlisting[language=C, style=mystyle,
%     caption={Código completo. \propia{}},
%     label={lst:codigo}]{programa/mi_programa.c}
% }{
%   \noindent\textbf{Nota:} El archivo aún no está disponible.
% }
%
%
% --- FORMA 3: Código inline (fragmentos pequeños) ---
% \begin{lstlisting}[language=Python, style=mystyle,
%   caption={Fragmento de ejemplo},
%   label={lst:ejemplo}]
% def fibonacci(n):
%     if n <= 1:
%         return n
%     return fibonacci(n - 1) + fibonacci(n - 2)
% \end{lstlisting}
%
%
% --- FORMA 4: Makefile ---
% \lstinputlisting[language=make, style=mystyle,
%   caption={Makefile del proyecto. \propia{}},
%   label={lst:makefile}]{codigos/Makefile}
%
%
% --- FORMA 5: Archivos de cabecera (.h) ---
% \lstinputlisting[language=C, style=mystyle,
%   caption={Archivo de cabecera. \propia{}},
%   label={lst:header}]{programa/common.h}
%
%
% ==========================================
% BUENAS PRÁCTICAS
% ==========================================
%
% 1. Usa \lstinputlisting para código que compila.
%    Así el PDF siempre refleja el código real.
%
% 2. Usa \IfFileExists para que la compilación no
%    falle si falta un archivo.
%
% 3. Pon siempre caption y label para referenciar
%    los listados desde el texto:
%    "En el Listado~\ref{lst:etiqueta} se muestra..."
%
% 4. Usa \propia{} al final del caption para agregar
%    "Fuente: Elaboración propia."
%
% 5. Para proyectos con varios capítulos, crea un
%    anexos.tex por capítulo y agrégalo arriba con
%    \input{}.


% ==========================================
\end{document}
% ==========================================
